\subsection{Further examples of affine Springer fibers}\label{ss:aff examples}
In this subsection we give more examples illustrating the rich geometry of affine Springer fibers. We omit the calculations that lead to the geometric descriptions. 

In all examples below, the affine Springer fibers are {\em homogeneous} in the sense that $\sX_{\gamma}$ is equipped an extra $\Gm$-action coming from the {\em loop rotation} on $\Gr_{G}$ by dilation on the uniformizer $t$. For more information on homogeneous affine Springer fibers and their application to representation theory, see \cite{OY}. 


\subsubsection{}\label{SL2 ram} When $G=\SL_{2}$ and $\gamma=\mat{0}{t^{m+1}}{t^{m}}{0}$.  When $m$ is even, $\sX_{\gamma}$ consists exactly of those lattices $\L$ such that
\begin{equation*}
t^{m/2}\calO_{F}\oplus t^{m/2}\calO_{F}\subset\L\subset t^{-m/2}\calO_{F}\oplus t^{-m/2}\calO_{F}.\end{equation*}
In this case, $\sX_{\gamma}$ coincides with the closure of the $L^{+}G$-orbit in $\Gr_{G}$ corresponding to the coweight $m\alpha^{\vee}/2$. 

When $m$ is odd, $\sX_{\gamma}$ consists exactly of those lattices $\L$ such that
\begin{equation*}
t^{(m+1)/2}\calO_{F}\oplus t^{(m-1)/2}\calO_{F}\subset\L\subset t^{-(m-1)/2}\calO_{F}\oplus t^{-(m+1)/2}\calO_{F}.
\end{equation*}
In this case, consider instead the affine Grassmannian $\Gr_{G'}$ of $G'=\PGL_{2}$, which contains $\Gr_{G}$ as a component. Then $\sX_{\gamma}$ can be identified with the closure of the $L^{+}G'$-orbit in $\Gr_{G'}$ corresponding to the coweight $m\alpha^{\vee}/2$. 

In either case, we have $\dim\sX_{\gamma}=m$.



\subsubsection{The Lusztig-Smelt examples}\label{sss:LS} Let $G=\GL_{n}$ and $\gamma\in\frg(F)$ with characteristic polynomial $P(x)=x^{n}-t^{m}=0$, where $(m,n)=1$. If a lattice $\L\subset F^{n}$ is stable under $\gamma$, it carries an action of the ring $R=\calO_{F}[X]/(X^{n}-t^{m})$. Let $K=\Frac(R)$, then $K= k\lr{s}$ with $x=s^{m}$ and $t=s^{n}$. Then the integral closure of $R$ in $K$ is $\tilR:=k\tl{s}$. The action of $\gamma$ on $F^{n}$ makes it a one-dimensional $K$-vector space. We fix a $K$-linear isomorphism $F^{n}\cong K$, under which a lattice in $F^{n}$ stable under $R$ is simply a {\em fractional $R$-ideal}, i.e., a finitely generated $R$-submodule of $K=\Frac(R)$.  Then $\sX_{\gamma}$ can be identified with the moduli space of fractional $R$-ideals.

The centralizer $G_{\gamma}(F)$ is simply $K^{\times}$, which acts on the set of fractional $R$-ideals by multiplication. This action clearly factors through $K^{\times}/R^{\times}$, which is the group of $k$-points of the local Picard group scheme $P_{\gamma}$.

There is an action of $\Gm$ on $K$ (by field automorphisms) given by scaling $s$ (so $s^{i}$ gets weight $i$ under this action). This induces an action of $\Gm$ on $\sX_{\gamma}$. The fixed points of $\Gm$ on $\sX_{\gamma}$ correspond to fractional ideals generated by monomials of $s$. More precisely, if a fractional $R$-ideal $\L\subset K=k\lr{s}$ is fixed by $\Gm$,  define $M_{\L}=\{i\in\ZZ|s^{i}\in \L\}$ which is a subset of $\ZZ$ stable under adding $m$ and $n$, because $\L$ is an $R=k\tl{s^{m},s^{n}}$-module. Therefore $M_{\L}\subset\ZZ$ is a finitely generated module for the monoid $A_{m,n}:=\ZZ_{\geq 0}m+\ZZ_{\geq0}n\subset\ZZ_{\geq 0}$. The assignment $\L\mapsto M_{\L}$ gives a bijection
\begin{equation*}
\sX_{\gamma}^{\Gm}(k)\isom\{\mbox{$A_{m,n}$-submodules $M\subset \ZZ$}\}.
\end{equation*}
Any $A_{m,n}$-submodule of $\ZZ$ contains all sufficiently large integers. Therefore any two such $A_{m,n}$-module $M$ and $M'$ differ by finitely many elements, and we can define $[M:M']=\#(M\bs M')-\#(M'\bs M)$. Fox any fixed $i\in\ZZ$, we have a total of $\frac{1}{n+m}\binom{n+m}{n}$ fixed points with $[M:\ZZ_{\geq0}]=i$. For a fixed point $p_{M}$ corresponding to an $A_{m,n}$-module $M$, consider the subvariety $C_{M}=\{p\in\sX_{\gamma}|\lim_{\Gm\ni z\to0}z\cdot p=p_{M}\}$. Then $C_{M}$ is isomorphic to an affine space whose dimension can be expressed combinatorially in terms of $M$. The cells $C_{M}$ give a stratification of $\sX_{\gamma}$. This gives a way to compute the Poincar\'e polynomial of connected components of $\sX_{\gamma}$. For more details, and the similar picture for $\sY_{\gamma}$, see \cite{LS}.

\subsubsection{} We look at the geometry of $\sY_{\gamma}$ in a special case of \S\ref{sss:LS}. We consider the case $G=\GL_{3}$ and $\gamma^{3}-t^{2}=0$. Introducing the variable $s$ with $t=s^{3}$ and $\gamma=s^{2}$ as before, then $R=k\tl{s^{2},s^{3}}$ with fraction field $K=k\lr{s}$. The affine Springer fiber  $\sY_{\gamma}$ classifies a chain of fractional $R$-ideals $\L_{0}\subset \L_{1}\subset \L_{2}\subset s^{-3}\L_{0}$.  We consider a component of $\sY^{0}_{\gamma}\subset\sY_{\gamma}$, classifying chains $\{\L_{i}\}$ as above with $[\L_{i}:k\tl{s}]=i$, $0\leq i\leq 2$. 

We first study the $\Gm$-fixed points on $\sY^{0}_{\gamma}$. For each $\Gm$-fixed $R$-submodule of $k\lr{s}$, we denote its associated module for the monoid $A_{3,2}$ (see \S\ref{sss:LS}) by a sequence of integers. For example, $(0,1,2,\cdots)$ stands for the standard lattice $k\tl{s}$. Note that the sequence for any $\Gm$-fixed $R$-fractional ideal is either consecutive $(n,n+1,n+2,\cdots)$ or has at most one gap at the second place $(n,n+2,n+3,\cdots)$. We have the following four fixed points in $\sY^{0}_{\gamma}$:
\begin{enumerate}
\item $q: (0,1,2,\cdots)\subset (-1,0,1,\cdots)\subset (-2,-1,0,\cdots)$;
\item $p_{0}: (-1,1,2,\cdots)\subset (-1,0,1,\cdots)\subset (-2,-1,0,\cdots)$;
\item $p_{1}: (0,1,2,\cdots)\subset (-2,0,1,\cdots)\subset (-2,-1,0,\cdots)$;
\item $p_{2}:(0,1,2,\cdots)\subset (-1,0,1,\cdots)\subset (-3,-1,0,\cdots)$. 
\end{enumerate}
Our indexing scheme is that $p_{i}$ is obtained from $q$ by changing the lattice $\L_{i}$. 

Then $\sY^{0}_{\gamma}$ is the union of three irreducible components $C_{0}\cup C_{1}\cup C_{2}$, and each component $C_{i}$ is isomorphic to $\PP^{1}$. They all contain $q$ and that is the only intersection between any two of them. We have $p_{i}\in C_{i}$ for $i=0,1,2$.

There is a natural way to index affine partial flag varieties  of $G$ by subsets $J\subset \{0,1,2\}$, as we saw in \S\ref{sss:GLn parahoric}. Let $\sX_{J,\gamma}$ be the affine Springer fiber of $\gamma$ in $\Fl_{J}$. Under the projection $\sY_{\gamma}\to \sX_{J,\gamma}$, the curves $C_{i}$ for $i\notin J$ collapse to a point, and the other curves map isomorphically onto their images.



\subsubsection{}\label{ex:Sp} Let $G=\Sp(V)$ where $V$ is a symplectic space of dimension $2n$ over $k$, and assume $\chark\neq2$. Fix a decomposition $V=U\oplus U^{*}$ into Lagrangian subspaces of $V$, such that the symplectic form restricts to the natural pairing on $U\times U^{*}$. Consider $\gamma=\mat{0}{tX}{Y}{0}\in\frg(F)$ where $X\in\Sym^{2}(U)$ (viewed as a self-adjoint map $\xi:U^{*}\to U$) and $Y\in\Sym^{2}(U^{*})$ (viewed as a self-adjoint map $\eta:U\to U^{*}$). The condition that $\gamma$ is regular semisimple is equivalent to that: (1) both $\xi$ and $\eta$ are isomorphisms; (2) $\xi\eta\in\GL(U)$ is regular semisimple, or equivalently $\eta\xi\in\GL(U^{*})$ is regular semisimple.

The affine Springer fiber $\sX_{\gamma}$ classifies self-dual lattices $\L\subset V\otimes_{k}F$ which are stable under $\gamma$ (see \S\ref{sss:ASF Gr GLn}). Consider the action of $\Gm$ on $V\otimes_{k}F=U\otimes_{k}F\oplus U^{*}\otimes_{k}F$ such that $Ut^{i}$ has weight $2i-1$ and $U^{*}t^{i}$ has weight $2i$. This induces a $\Gm$-action on $\Gr_{G}$ and on $\sX_{\gamma}$. We first consider the fixed points $\Gr_{G}^{\Gm}$. A lattice $\L\in\Gr_{G}$ is fixed under $\Gm$ if and only if it is the $t$-adic completion of
\begin{equation*}
\Span\{(A_{i}\oplus B_{i})t^{i}; i\in\ZZ \}
\end{equation*}
where $\cdots \subset A_{-1}\subset A_{0}\subset A_{1}\subset \cdots \subset U$ is a filtration of $U$ such that $A_{N}=0$ for $N\ll 0$ and $A_{N}=U$ for $N\gg 0$, and $\cdots \subset B_{-1}\subset B_{0}\subset B_{1}\subset \cdots \subset U^{*}$ is a similar filtration of $U^{*}$, such that
\begin{equation}\label{AB dual}
B_{i-1}=A^{\bot}_{-i}, \forall i\in\ZZ
\end{equation}
under the duality pairing between $U$ and $U^{*}$. The last condition reflects the fact that $\L$ is self-dual  under the symplectic form.

A lattice $\L\in\sX_{\gamma}^{\Gm}$ is then determined by two filtrations $A_{\bu}$ of $U$ and $B_{\bu}$ of $U^{*}$, dual in the sense \eqref{AB dual}, with the extra condition that
\begin{equation*}
\eta A_{i}\subset B_{i}; \quad \xi B_{i}\subset A_{i+1}, \quad\forall i\in\ZZ.
\end{equation*}
We summarize the data into the following diagram
\begin{equation*}
\xymatrix{ A_{-1}\ar@{^{(}->}[rr]\ar[dr]^{\eta} && A_{0}\ar@{^{(}->}[rr]\ar[dr]^{\eta} && A_{1}\ar[dr]^{\eta}& \cdots\\
\cdots & A^{\bot}_{0}\ar@{^{(}->}[rr]\ar[ur]^{\xi} && A^{\bot}_{-1}\ar@{^{(}->}[rr]\ar[ur]^{\xi} && A_{-2}^{\bot}}
\end{equation*}

For example, when $\dim V=4$, hence $\dim U=2$, there are two possibilities. The first possibility is $A_{-1}=0\subset A_{0}=U$, which corresponds to the standard lattice $V\otimes_{k}\calO_{F}$. The second possibility is $A_{-1}=0\subset A_{0}\subset A_{1}=U$, and $A_{0}\subset U$ is a line satisfying $\xi A^{\bot}_{0}\subset A_{0}$, i.e., $A^{\bot}_{0}\subset U^{*}$ is an isotropic line under the quadratic form $X$. There are two such lines $A^{\bot}_{0}$, giving two other $\Gm$-fixed points. Therefore, $\sX_{\gamma}^{\Gm}$ consists of 3 points. We have $\sX_{\gamma}=C_{1}\cup C_{2}$ where $C_{i}\cong\PP^{1}$, the two components intersect at the standard lattice, and each $C_{i}$ contains one of the remaining $\Gm$-fixed points.


\subsubsection{The Bernstein-Kazhdan example}\label{sss:BK example} In \cite[Appendix]{KL}, Bernstein and Kazhdan gave the first example of an irreducible component of an affine Springer fiber which was not a rational variety. We keep the same notation as in \S\ref{ex:Sp}. Let $\Fl_{\bP}$ be the partial affine flag variety of $G=\Sp(V)$ classifying pairs of lattices $\L'\subset\L$ such that $\L^{\vee}=\L$ and $\L'^{\vee}=t^{-1}\L$.
Let $\gamma$ be as in \S\ref{ex:Sp}. Then the same $\Gm$ acts on $\sX_{\bP,\gamma}$, and the fixed points $\L'\subset \L$ can be described by two pairs of filtrations $(A_{\bu}, B_{\bu})$ and $(A'_{\bu}, B'_{\bu})$, where $(A_{\bu}, B_{\bu})$ is the kind of filtration of $U$ and $U^{*}$ as described in \S\ref{ex:Sp}, and $(A'_{\bu}, B'_{\bu})$ is similar except that \eqref{AB dual} is replaced by $B'_{i}=A_{-i}'^{\bot}$. Moreover, the inclusion $\L'\subset \L$ is equivalent to $A'_{i}\subset A_{i}$ and $B'_{i}\subset B_{i}$, for all $i$.

Consider for example $\dim V=6$ and we fix the dimension of the filtrations:
\begin{equation*}
A_{-1}=0, \quad \dim A_{0}=2, \quad A_{1}=U;\quad  \dim A'_{0}=1, \quad A'_{1}=U.
\end{equation*}
Such filtrations $(A_{\bu}, B_{\bu}; A'_{\bu}, B'_{\bu})$ are determined by the complete flag $0\subset A'_{0}\subset A_{0}\subset U$ satisfying $\xi A^{\bot}_{0}\subset A_{0}$ and that $\eta A'_{0}\subset A'^{\bot}_{0}$. In other words,  $A^{\bot}_{0}$ is an isotropic line in $U^{*}$ under the quadratic form $X$, and $A'_{0}$ is an isotropic line in $U$ under the quadratic form $Y$. The pair $(A^{\bot}_{0}, A'_{0})$ determines a point in $Q(X)\times Q(Y)\subset \PP(U)\times\PP(U^{*})$, the product of conics defined by $X$ and $Y$. The incidence relation $A'_{0}\subset A_{0}$ defines a curve of bidegree $(2,2)$ in $Q(X)\times Q(Y)\cong\PP^{1}\times \PP^{1}$, which is then a curve of genus one. Therefore, a connected component of $\sX^{\Gm}_{\bP,\gamma}$ is a curve of genus  one. Consider the points in $\sX_{\bP,\gamma}$ that contract to this curve, and take its closure $Z$. One can show that $\dim Z=3=\dim\sX_{\bP,\gamma}$. Hence $\sX_{\bP,\gamma}$ contains an irreducible component $Z$ which is irrational. We refer to the appendix of \cite{KL} for more details.



\subsubsection{``Subregular'' affine Springer fibers}
When $\sX_{\gamma}$ or $\sY_{\gamma}$ is one-dimensional, we may call them subregular affine Springer fibers, by analog with subregular Springer fibers discussed in \S\ref{sss:subreg}. If $\dim\sY_{\gamma}=1$, it is a union of $\PP^{1}$'s, hence we can define its dual graph. In \cite[Prop 7.7]{KL}, the dual graphs of the subregular affine Springer fibers in $\Fl$ are classified, and they are almost always the extended Dynkin diagrams of simply-laced groups, except that they can also be infinite chains in type $A$ (see Example \ref{P1 chain}). 



